\subsection{Related Works}

\smallsec{LLMs for Theorem Proving in Lean} Formal theorem proving is a rapidly growing frontier in AI for mathematics and software verification \citep{yang2024formal, li2024survey}. Progress is typically measured with benchmarks of mathematical theorems in Lean such as miniF2F \citep{zheng2021minif2f}, PutnamBench \citep{tsoukalas2024putnambench}, and ProofNet \citep{azerbayev2023proofnet}. Recently, there have been many LLMs developed for Lean such as Seed-Prover \citep{chen2507seed}, Goedel-Prover \citep{lin2025goedel}, DeepSeek-Prover \citep{ren2025deepseek}, and Kimina-Prover \citep{wang2025kimina}.  There have also been post-training techniques built on top of these models, such as with expert iteration \citep{lin2024lean}, proof sketching \citep{cao2025reviving}, tree search \citep{lample2022hypertree, zimmer2025bourbaki}, self-play \citep{dong2025stp}, proof repair \citep{ospanov2025apollo}, and RL \citep{gloeckle2024abel}.

\smallsec{AI for Program Simplification} A related line of work makes programs shorter or more efficient \citep{schkufza2013stochastic, alphadev, pie, gautam2024refactorbench}. In parallel, library learning aims to discover reusable abstractions, often eliminated repeated code and shortening programs \citep{ellis2023dreamcoder, grand2023lilo, kaliszyk2015learning, wang2023lego, zhou2024refactor, berlot2024library}. Finally, symbolic reasoning techniques like program slicing \citep{weiser2009program}, super-optimization \citep{sasnauskas2017souper}, or partial evaluation \citep{jones1996introduction} can also shorten and optimize low-level code. 

\smallsec{Automated Proof Shortening} ~\citet{frieder2024data} study factors that make Lean proofs easier to understand, motivating shorter proofs for maintainability. Classically, there have also been many symbolic methods targeting shortening proofs in SAT and first-order logic languages ~\citep{rahul2001proof, vyskovcil2010automated, wernhard2024investigations, gladshtein2024small,kinyon-proof-simp}. On the neural side, GPT-f~\citep{polu2020generative} generated 23 verified proofs shorter than those in the Metamath library. Most related to our work, ImProver~\citep{ahuja2024improver}, is an inference-time method for proof shortening using GPT-4o with proof states and retrieval. In contrast, we use training-time approaches (expert iteration and RL), analyze complementary inference-time techniques, and focus on shortening longer proofs generated by SoTA LLMs.